% RESEARCH_DRAFT_SHA256: 67d46fd85eb443fe500d17a0dae03ee5cbc1d9e508091d23a929499bdd926dfd
% Layout and source-reference wrappers expanded for the web renderer.
\documentclass[11pt]{article}
\usepackage[margin=25mm]{geometry}
\usepackage{fontspec,unicode-math}
\setmainfont{Latin Modern Roman}
\setmathfont{Latin Modern Math}
\usepackage{amsmath,amssymb,mathtools,xurl,hyperref}
\hypersetup{colorlinks=true,urlcolor=blue,linkcolor=blue}
\setlength{\parindent}{0pt}
\setlength{\parskip}{5pt}
\setlength{\emergencystretch}{3em}
\newcommand{\sref}[2]{\hyperlink{source:#1:#2}{[S#2]}}
\newcommand{\cataloguescope}[1]{\paragraph{Source question.}#1}
\newcommand{\floor}[1]{\left\lfloor#1\right\rfloor}
\begin{document}

\section*{Operational distances for quantum functions and an operator-norm testing obstruction}
\noindent catalogue entry; catalogue entry. Research attempt dated 11 October 2026.
\paragraph{Result in scope.}This manuscript contains a source-specific attempt and complete proofs of the stated auxiliary calculations. It makes no new solution or novelty claim for the entire target.
\section{Original problem and preserved definition}
The following definition and source list are reproduced from the repository. Any earlier inspection dates and reported literature claims in this preserved text describe that source record, not additional verification by this attempt.
% BEGIN PRESERVED DEFINITION
\subsection{Definitions and mathematical statement}

\paragraph{Shared notation.}Standard number systems are denoted by $\mathbb N,\mathbb Z,\mathbb Q,\mathbb R,\mathbb C$; index ranges, computational models and specialized notation follow the cited problem source and the definitions below.


In the source model a finite-dimensional quantum function $F:X\to A$ assigns positive-semidefinite $d\times d$ measurement operators $M_{a\mid x}$ satisfying $\sum_{a\in A}M_{a\mid x}^2=I$ for each $x$. Sequential query probabilities are $d^{-1}\|M_{a_1\mid x_1}\cdots M_{a_q\mid x_q}\|_{\mathrm F}^2$, with the order interpreted as in the source. The function is commutative if all its measurement operators commute. A property is a class of such functions; a property tester seeks to distinguish membership from being far from the class using few queries.
\paragraph{Statement recorded in the supplied source.}
What kinds of noncommutative properties are there? How does one measure a quantum function’s distance to a noncommutative property? What properties of quantum functions can be tested using few queries? \href{https://doi.org/10.4171/owr/2021/54}{[S2]}
\paragraph{Scope qualification.}
Choosing the distance is explicitly part of this research question. The source does not supply a fixed metric or dimension-independent query bound. These operators need not be projections.
\paragraph{Source status.}
The supplied catalogue labels this record open. Its source-status metadata was checked against the supplied snapshot on 2026-10-11. This import does not independently establish current open status or mathematical correctness. The catalogue literature assessment is dated 2026-08-22.
\subsection{Short English statement}
Develop notions of properties, distance, and efficient testing for quantum functions with noncommuting measurements.
\subsection{Sources}
\begin{itemize}
\item[S1] Ulam AI and the UnsolvedMath Contributors, UnsolvedMath: A Curated Collection of Open Mathematics Problems, catalogue entry. Catalogue curation and metadata: CC BY 4.0. The supplied source question is retained; definitions and the English summary are editorial additions.
\url{https://huggingface.co/datasets/ulamai/UnsolvedMath}.
\item[S2] Complexity Theory (2022), Oberwolfach Reports, DOI 10.4171/owr/2021/54, pp. 2955--3005
\url{https://doi.org/10.4171/owr/2021/54}.
\item[S3] Complexity Theory, Oberwolfach Reports 54/2021, primary report PDF
\url{https://ems.press/content/serial-article-files/46932}.
\textit{Underlying Oberwolfach Reports material retains its source-specific licence, including CC BY-SA 4.0 where applicable; see THIRD\_PARTY\_NOTICES.md.}
\end{itemize}
% END PRESERVED DEFINITION
\subsection{Verbatim supplied catalogue question}
What kinds of noncommutative properties are there? How does one measure a quantum function’s distance to a noncommutative property? What properties of quantum functions can be tested using few queries?

\section{The source question and the attempted metric}
The primary report, printed pages 2972--2973, asks for properties,
distances and few-query tests, explicitly leaving the distance
unspecified. We try a transcript distance and check whether it
controls a natural operator-norm notion of commutativity. The
answer to that attempted implication is negative. We retain the
literal positive-measurement convention
$\sum_aM_{a\mid x}^2=I$; no projectivity is assumed except in
our concrete example. This does not characterize all testable
properties in the source's broad agenda.

\section{An exact operational pseudometric}
Assume finite input and output alphabets. For a deterministic
adaptive query tree $T$ of depth at most $q$, let $P_F^T$ be its
output-transcript distribution when querying $F$, and define
\[
 d_q(F,G)=\max_T\frac12\sum_\omega
 |P_F^T(\omega)-P_G^T(\omega)|.
\]
There are finitely many such trees. The probabilities are valid
because summing over the next measurement outcome uses
$\sum_aM_{a\mid x}^2=I$ and preserves the prefix probability.
The expression is nonnegative and symmetric, vanishes at identical
boxes, and obeys the triangle inequality by applying it termwise
for each tree and then taking maxima. Ignoring an extra query
shows $d_q\le d_{q+1}$. Thus
\[
 d_{\rm seq}(F,G)=\sum_{q\ge1}2^{-q}d_q(F,G)
\]
is a well-defined pseudometric bounded by one. It vanishes exactly
when all finite adaptive transcript laws agree, since all summands
are nonnegative. It need not distinguish operator presentations.
For example simultaneous unitary conjugation leaves every product
Frobenius norm unchanged. Replacing every $M$ by $M\otimes I_k$
also preserves the normalized squared norm, because both squared
Frobenius norm and Hilbert-space dimension multiply by $k$.

Any depth-$q$ tester's acceptance probabilities on two boxes differ
by at most $d_q(F,G)$. For a deterministic tree and an acceptance
function $a(\omega)\in[0,1]$, the signed difference is
$\sum_\omega a(\omega)(P_F^T-P_G^T)(\omega)$, whose absolute
value is at most the sum of its positive masses, namely total
variation. Randomized trees are convex combinations, so the same
bound persists. Conversely, for a maximizing tree, accepting
exactly transcripts with positive probability difference attains
total variation. This last statement is information-theoretic;
it does not say the best tree is efficiently computable.

\section{Small operational distance can hide noncommutativity}
For $D\ge1$ form a direct sum of a $D$-dimensional deterministic
box and a two-dimensional projective box with two binary queries.
For the first query the two-dimensional outcome-zero projector is
\[
 P=\begin{pmatrix}1&0\\0&0\end{pmatrix},
\]
and for the second it is
\[
 Q=\frac12\begin{pmatrix}1&1\\1&1\end{pmatrix}.
\]
The other projectors are $I-P,I-Q$. On the first block both
queries deterministically output zero. Products remain block
diagonal, so every transcript distribution of the combined box
$F_D$ is exactly the mixture
\[
 \frac{D}{D+2}P_{\rm det}^T+\frac2{D+2}P_{P,Q}^T.
\]
This formula also holds for adaptive trees: the same fixed block
is used along each transcript. Hence for the fully deterministic
commuting box $G_D$ of dimension $D+2$,
\[
 d_q(F_D,G_D)\le\frac2{D+2}\quad\hbox{for every }q.
\]
The bound concerns one sequential run of the box; a newly reset
independent copy would sample its block anew.

Nevertheless $F_D$ is a constant distance from commuting
operators in the following precise sense. Its two outcome-zero
operators are $\widetilde P=I_D\oplus P$ and
$\widetilde Q=I_D\oplus Q$, and
\[
 [P,Q]=\frac12\begin{pmatrix}0&1\\-1&0\end{pmatrix},
 \qquad\|[\widetilde P,\widetilde Q]\|_{\rm op}=1/2.
\]
If commuting positive contractions $P',Q'$ on the same space
both approximate these operators within $\delta$ in operator
norm, then
\[
 \|[\widetilde P,\widetilde Q]\|_{\rm op}
 \le\|[\widetilde P-P',\widetilde Q]\|_{\rm op}
    +\|[P',\widetilde Q-Q']\|_{\rm op}\le4\delta.
\]
Thus $\delta\ge1/8$. Positive normalized measurement operators
are contractions, since $M_a^2\preceq I$, so the conclusion
applies in particular to every commuting measurement family on
the same space. It even allows $P',Q'$ to mix the two blocks.

A uniform test that accepts commuting boxes with probability at
least $2/3$ and rejects all boxes at this operator-norm distance
with probability at least $2/3$ would need an acceptance gap at
least $1/3$. For $D>4$ that contradicts the transcript bound,
regardless of the length of a single sequential run. Even if a
fixed number $q$ of fresh-copy runs were allowed, coupling the
block choices gives a total variation bound at most
$2q/(D+2)$, by the union bound on ever choosing the small block.
Consequently no dimension-independent constant-copy guarantee
follows in that model either. This does not preclude tests in
normalized Hilbert--Schmidt distance or with a dimension-dependent
sample budget.

\section{Audit and the first unresolved step}
The attempted passage from operational closeness to uniform
operator closeness fails through a low-weight direct-sum sector.
The audit checked adaptive transcripts, dimension normalization,
operator contractions, possible mixing of sectors by the
approximant, and the distinction between one run and fresh copies.
A useful general property theory still needs a chosen equivalence
and distance, robust positive results, and efficient test
constructions. No numerical experiment or theorem prover was
used, and the elementary distance construction and obstruction
are not claimed as new original-target progress.

\section{Process and contributions}
Mathematics and model-assisted adversarial self-audit: OpenAI gpt-6-astra, ultra reasoning, as dispatched for this research batch; exact serving revision not exposed. Codex (GPT-6; exact serving revision not exposed) materially assisted with orchestration and format review. No human mathematical authorship, mathematical edits, or independent human proof checking is asserted.
The coordinating Codex (GPT-6; exact serving revision not exposed) also provided a separate model-assisted written-proof review of this manuscript. This is not independent human verification.
Contribution prompt A was applied to the full definition and supplied source row, then prompt B was applied adversarially to the derived results. The item-specific audit above records the scope checks and rejected extrapolations. Local assembly and structural checks do not establish mathematical correctness. The source catalogue status is preserved in TOP\_PROBLEM and is separate from this attempt status.
\paragraph{Integration review.} Codex (GPT-6; exact serving revision not exposed) reviewed the written arguments and expanded layout and source-reference wrappers. Repository metadata and the production web renderer were checked using Python 3.13 and JavaScript with MathJax 3.2.2. This model-assisted review is not independent human proof verification; no PDF compilation or formal proof checking is asserted.

\section{Completion Estimate}
\noindent\textbf{COMPLETION: 0\%}
The proved material is a rederivation of known or auxiliary facts rather than a substantiated new advance on the original target. The transcript pseudometric is one possible operational choice, and the obstruction rules out one stronger operator-norm formulation. A general characterization of noncommutative properties and efficient robust tests is not obtained.
\section{References}
The complete original Sources subsection above is retained, including catalogue attribution and licence notices. The additional list records primary sources actually consulted for this attempt on 11 October 2026; uncited source links reproduced above were not all rechecked.
\begin{itemize}
\item Henry Yuen, Noncommutative Property Testing, OWR54/2021 pp.2971–2974; full primary source passage read in cached text, and page2972 operator definition visually checked from the primary PDF. The original agenda, stated known results and convention mismatch are distinguished. \url{https://ems.press/content/serial-article-files/46932}.
\end{itemize}
\section{Final status}
\textbf{Unresolved in this attempt.} The transcript pseudometric is one possible operational choice, and the obstruction rules out one stronger operator-norm formulation. A general characterization of noncommutative properties and efficient robust tests is not obtained.
\end{document}
